

Le choix du langage à utiliser était libre, nous avons donc opté pour un langage simple d'utilisation et ayant déjà fait ses preuves : \textbf{Java}.

\begin{figure}[!h]
\centering
\includegraphics[width=8cm]{images/javaLogo.png}
\caption{Logo Java}
\label{logo_java}
\end{figure}

nous avons choisi de travailler avec le logiciel \textbf{Éclipse}, l'\gls{IDE} par excellence pour développer du \textbf{Java}.

\begin{figure}[H]
\centering
\includegraphics[width=8cm]{images/eclipseLogo.png}
\caption{Logo Éclipse}
\label{logo_eclipse}
\end{figure}

\textbf{Éclipse} est agréable à utiliser, il est libre de droits et fait partit des IDE les plus puissants de nos jours.

La bibliothèque nous permettant de faire des \gls{ihm} utilisé est swing, une bibliothèque native de Java, constamment mise à jours et très utilisé (donc très documenté).

Nous avons choisi \textbf{Git} pour gérer les versions de notre projet, actuellement hébergé sur un serveur \textbf{GitHub}.

\begin{figure}[H]
\centering
\includegraphics[width=8cm]{images/githubLogo.png}
\caption{Logo GitHub}
\label{logo_github}
\end{figure}

Pour finir, nous avons voulu organiser la répartition des taches avec l'outil en ligne : \textbf{Trello}.

\begin{figure}[H]
\centering
\includegraphics[width=8cm]{images/trelloLogo.png}
\caption{Logo Trello}
\label{logo_trello}
\end{figure}